\section{Table Types}\label{sec:table-types}

Each table in a table file has a table type corresponding to the value defined in the
table. For example, the \tabletype{thrust} table type defines a thrust value. The table
type is compared with its use in the problem file, and an input error occurs if the
table is not used correctly---for example, if a thrust table is used for an axial-force
coefficient. \Cref{tab:table-types} lists the table types used in TAOS.

\begin{table}[htbp]
  \centering
  \caption{Table types.}
  \label{tab:table-types}
  \begin{tabularx}{0.82\linewidth}{@{}>{\ttfamily}lX@{}}
    \toprule
    \normalfont Table type & Description \\
    \midrule
    ca     & Axial-force coefficient (in body coordinates) \\
    cn     & Normal-force coefficient (in body coordinates) \\
    cl     & Lift-force coefficient (in wind coordinates) \\
    cd     & Drag-force coefficient (in wind coordinates) \\
    cs     & Side-force coefficient (in wind coordinates) \\
    cx     & Body $x$-axis force coefficient \\
    cy     & Body $y$-axis force coefficient \\
    cz     & Body $z$-axis force coefficient \\
    thrust & Thrust magnitude \\
    tvec1  & Thrust-vector angle 1 \\
    tvec2  & Thrust-vector angle 2 \\
    mdot   & Mass flow \\
    cg     & Center of gravity \\
    windv  & Wind magnitude \\
    windh  & Wind direction or heading \\
    winde  & Wind component in east direction \\
    windn  & Wind component in north direction \\
    windd  & Wind component in ``down'' direction \\
    output & User-defined output variable \\
    \bottomrule
  \end{tabularx}
\end{table}

Aerodynamic-force coefficients can be defined in one of three ways: in the vehicle
body axes with $C_X$, $C_Y$, and $C_Z$; in the wind axes with $C_L$, $C_D$, and
$C_S$; or in an axisymmetric system with $C_A$ and $C_N$. These sets of
coefficients go together and cannot be mixed. For example, $C_L$, $C_D$, and
$C_S$ form a valid set, whereas $C_L$, $C_D$, and $C_A$ do not.

Thrust and mass-flow tables are used for all types of jet and rocket motors. The
thrust table defines the magnitude of the thrust vector; the direction of the thrust
vector relative to the body coordinate system is given by the thrust-vector-angle
tables.

Point-mass trajectory-analysis methods assume that a vehicle is always in a trimmed
condition; that is, the moments always sum to zero. Therefore, trimmed aerodynamic
coefficients should be input rather than untrimmed values. Trimmed aerodynamic
coefficients are often a function of the vehicle's center of gravity, and center of
gravity is in turn often a function of vehicle weight. The \tabletype{cg} table type is
provided to define a center-of-gravity value that can be used later in the aerodynamic
coefficient tables.

Wind tables define the components of the wind-velocity vector in a local-horizon
coordinate system. Horizontal wind velocity is given either by the \tabletype{windv}
(velocity magnitude) and \tabletype{windh} (wind heading) tables or by the
\tabletype{winde} (east component) and \tabletype{windn} (north component)
tables. The \tabletype{windv} and \tabletype{windh} tables go together, and the
\tabletype{winde} and \tabletype{windn} tables go together. A vertical wind
component can also be input with the \tabletype{windd} table.

Output tables are used to calculate user-defined output variables such as heat-transfer
rates. Values from output tables are not used during the trajectory calculations, but
they are available when computing user-defined output variables
(\cref{sec:trajectory-block-define,sec:problem-block-define}).
